\section{The remaining multiplicity in the denominator argument}

% BEGIN F86D1
\paragraph{F86D1. One numerical test per localization.}
Fix $i$. For each good point $x$ choose a witnessing list $\Sigma_x$
and write $A_x=M_{j,\Sigma_x}$ and $E_x=\mathrm{eigC}(i,x)$.
The subspace $A_x$ is independent of the choice. In fact any two
witnesses have upper labels over the same point. Their augmented
base-change lists agree up to permutation by H86P3. Undo their common
$\eta^{-1}$ twist and subtract one copy of the marked length-$t$
term from the two multisets. The remaining lists agree up to
permutation; this preserves both the local sums and the inequalities
defining a row parameter. The same lower summands therefore generate
the two spaces.

By J86O3 the upward map $T$ takes $A_x$ into $E_x$.
The $E_x$ at different points are mutually orthogonal, since they are
sums of distinct upper types at this level, by U86T3. By V86E1 the
good row space is the sum of the $A_x$. On the entire row space $T$
is injective by L86M2. Thus
\[
 d_M(i)=\dim T(\mathrm{goodPartM}(i))
       =\sum_{x\ \mathrm{good}}\dim TA_x
       \ \le\ \sum_{x\ \mathrm{good}}\dim E_x
       =\sum_{x\ \mathrm{good}}\mathrm{mult}(i,x).
\]
All these sums are finite. Their nonnegative integer difference
is $\sum_x\dim(E_x/TA_x)$. Consequently the single reverse inequality
\[
             \sum_{x\ \mathrm{good}}\mathrm{mult}(i,x)\le d_M(i)
             \tag{2}
\]
is equivalent to $TA_x=E_x$ for every good point and every list over
it. By V86E2 this is equivalent to the joint downward membership and
downward injectivity conditions there. It also gives the original
$\mathrm{ConclusionZ}$ at $i$: the restricted kernels vanish by
injectivity of $T$, equality of the images gives the good-part
identity after taking the sum, and (2) is equality by the displayed
calculation. Conversely $\mathrm{ConclusionZ}$ supplies (2) at
nonempty good support by its last clause. On empty good support
V86E1 gives $\mathrm{goodPartM}(i)=0$, so both sides of (2) vanish
and the pointwise tests are vacuous. Thus the guard $\mathrm{HG}$
does not require an extra existence hypothesis for this equivalence.
% END F86D1

% BEGIN F86D2
\paragraph{F86D2. The remaining original member tests.}
In $\mathrm{PointConditions}'$, S86A1--4 give all ten entries of
$\mathrm{RowsAAlt}$, U86T3 gives $\mathrm{RowsB}$, and L86M2 gives
the final rank comparison for every size datum. In $\mathrm{RowsC}$,
J86O2 gives joint occurrence on the specified residual representation,
R86G7 gives the integral rank-$r$ representation, R86G4 gives the
full degree of the tame polynomial, J86O3 gives the upward containment
and V86E1 identifies the good row space. Its downward containment entry follows by the last part of V86E2
and the spanning in F86I2. Paired with the separate $\mathrm{typeNZ}$
field of $\mathrm{PointConditions}'$, it gives the two downward tests. F86D1 proves that (2) at every
$i$ is both necessary and sufficient for these two entries together.
Thus the two original fields $\mathrm{what}_2$ and $\mathrm z$ are
each equivalent on this member to (2) at all its localizations.

For $\mathrm{AdmissibleDatum}$, S86A2--4, H86P1, U86T1 and I86A3
give the local, base-change, Galois and split theta conditions.
The original large coefficient duty in M84T6 supplies
$\mathrm{hypO}$. S86A1 supplies $\mathrm{adm5}$.
The residual irreducibility condition $\mathrm{adm6}$ is the very
predicate proved in R86G6, by unfolding its definition. The last
condition is $\mathrm{adm8}$: the indicated finite types must form
an internal direct sum in every allowed lower level space.
F86I1--2 prove the missing spanning comparison for every
$1\le j\le m$ using the reference-level multiplicity and the
generic packet on the two inner forms. Together with L86M5 this
gives the internal direct sum, so the entire $\mathrm{adm}$ holds.

Sixteen of the eighteen original preconditions are now proved.
U86T4 provides
$\mathrm{Conclusion1}$, R86G6 and R86G3 give $\mathrm{rhoBar}$ and
$\mathrm{g3p}$, L86M4 supplies $\mathrm{lam}$, and J86S3 supplies
$\mathrm{trip}$. The other ten are the original fields already
established through edition 0.85. The four ties have simultaneous
witnesses by M84T10, T86C6, M84T2 and T86C4 in their defining order.
Hence (2) on the selected members supplies all the member
inputs in A85D3. Its other twist, published, finiteness, summability
and uniform family hypotheses are retained. The remaining reverse inequality (2) has been used only as a hypothesis.
% END F86D2
